\documentclass[11pt]{article}

\usepackage[a4paper,margin=25mm]{geometry}
\usepackage[T1]{fontenc}
\usepackage{lmodern}
\usepackage{microtype}
\usepackage{amsmath,amssymb,mathtools}
\usepackage{booktabs,tabularx,array}
\usepackage{enumitem}
\usepackage{xcolor}
\usepackage[hidelinks]{hyperref}

\definecolor{prospectivered}{RGB}{185,0,0}
\newenvironment{prospective}{\color{prospectivered}}{}
\newcolumntype{Y}{>{\raggedright\arraybackslash}X}
\setlist{nosep,leftmargin=*}
\setlength{\emergencystretch}{2em}

\title{Common-Space Construction of Periodic Defect Operators}
\author{Eugene Joseph M. Ragasa\\
Department of Physics, De La Salle University, Manila, Philippines\\
\texttt{eugene.ragasa@dlsu.edu.ph}}
\date{Unpublished pre-results manuscript}

\begin{document}
\maketitle

\begin{abstract}
A defect-induced operator is not defined by subtracting independently constructed
Hamiltonian matrices merely because their dimensions agree. The pristine and defect
operators may live on different lattices, Bloch bundles, projected subspaces, localized
frames, orbital orderings, and energy references. This paper develops a fail-closed
common-space construction for periodic defect operators. The two calculations are first
placed on a matched comparison lattice and Brillouin torus. Constant rank, projector
compatibility, reciprocal-boundary sewing, and bundle topology are then screened before
an overlap-derived unitary is admitted. Wannier centers, site and orbital semantics, and
surviving symmetries constrain the remaining correspondence. A finite-dimensional
grammar supplies exact synthetic controls with prescribed rank, principal angles, polar
alignment, scalar offsets, isolation, and operator difference. A separately qualified
scalar energy reference completes the alignment, after which the pristine operator can
be transported to the defect space and subtracted. The resulting impurity operator is
gauge covariant rather than entrywise gauge invariant. Rank-changing defect states are
handled by an explicit common-core projection or enlarged retained space, never by zero
padding. The protocol separates geometry, projection, localization, alignment, energy,
numerical, and finite-supercell errors and assigns explicit blocked or unresolved
dispositions when a common space cannot be certified. Phosphorus and boron in silicon
are stated only as future specialization contracts. No production defect calculation,
retained common-space map, or material impurity operator is reported.
\end{abstract}

\noindent\textbf{Keywords:} defect Hamiltonian; common state space; operator alignment;
principal angles; polar decomposition; Wannier representation; impurity operator

\section{Introduction}
The difference between two linear operators is meaningful only when their domains and
codomains have been identified. In first-principles defect calculations this elementary
requirement is obscured by matrix output: pristine and doped workflows can produce
arrays of the same shape even though their supercell coordinates, target projectors,
Wannier gauges, orbital meanings, and additive energy zeros differ. Direct subtraction
then mixes a physical perturbation with representation mismatch.

Localized Wannier representations provide a natural finite periodic language for
operator comparison~\cite{marzari2012,Mostofi2008,Pizzi2020}. They do not, however, make
independently selected retained spaces identical. Equal rank provides an abstract
pointwise unitary map, not a unique physical correspondence. Principal angles and polar
decomposition quantify and align subspaces~\cite{bjorck1973,higham1986}, but they must be
embedded in a contract that also controls lattice folding, periodic sewing, orbital
semantics, surviving symmetry, and energy alignment.

For defect label $d$, the intended object is
\begin{equation}
 \Delta H_d^{(P)}
 =\overline H_d^{(P)}-U_d\overline H_b^{(P)}U_d^\dagger,
 \label{eq:target}
\end{equation}
where $H_b^{(P)}$ and $H_d^{(P)}$ are pristine and defect projected operators, $U_d$
identifies their retained spaces, and bars denote a shared scalar energy convention.
Every element of Eq.~\eqref{eq:target} is conditional on a declared physical and
numerical construction. The formula is the endpoint of alignment, not its substitute.

This paper formalizes that endpoint and its admission gates. It treats the common-space
problem independently of the scalar-relativistic bulk-silicon reduction of Paper
4~\cite{Ragasa2026d}, the general obstruction theory reserved for Paper
6~\cite{Ragasa2026f}, and the SOC impurity application reserved for Paper
7~\cite{Ragasa2026g}. The result is a prospective methodology. No numerical threshold,
material specialization, or execution evidence is implied by the definitions below.

\begin{prospective}
\section{Comparison geometry and retained spaces}

\subsection{Matched periodic geometry}
Let $\mathcal L_c$ be the common comparison lattice and $\mathbb T_c^D$ its Brillouin
torus. For a periodic defect supercell, the default comparison is a pristine calculation
repeated in the same supercell:
\begin{equation}
 H_b(\mathbf k),H_d(\mathbf k):\mathcal H_{\mathrm{amb},\mathbf k}
 \longrightarrow\mathcal H_{\mathrm{amb},\mathbf k},
 \qquad \mathbf k\in\mathbb T_c^D.
\end{equation}
The two branches must share lattice vectors, reciprocal-coordinate convention,
plane-wave or other ambient representation, boundary convention, units, relativistic
and spin convention, and a declared coordinate ordering. Atomic positions may differ
because of defect relaxation, but the mapping between comparison sites and translations
must remain explicit.

A primitive-cell pristine calculation may be folded into the defect supercell only
through a recorded isometry that specifies reciprocal preimages, degeneracy ordering,
normalization, and translation phases. Folded and directly recalculated pristine
operators must not be interchanged silently. Agreement of their selected invariants is a
validation test, not an identity assumed from equal spectra.

\subsection{Projected fibers}
For $s\in\{b,d\}$, let
\begin{equation}
 P_s(\mathbf k)^2=P_s(\mathbf k)=P_s(\mathbf k)^\dagger,
 \qquad \mathcal E_{s,\mathbf k}=\operatorname{Range}P_s(\mathbf k),
\end{equation}
and define
\begin{equation}
 H_s^{(P)}(\mathbf k)=P_s(\mathbf k)H_s(\mathbf k)P_s(\mathbf k).
\end{equation}
The projector construction records target bands, outer and frozen windows,
disentanglement, excluded states, and any defect-localized state policy. A common matrix
size after localization does not retroactively establish compatible projectors.

The primary route requires constant equal rank
\begin{equation}
 \operatorname{rank}P_b(\mathbf k)=\operatorname{rank}P_d(\mathbf k)=M
\end{equation}
throughout $\mathbb T_c^D$. The spectral isolation or entangled-subspace stability used
to define each projector is checked separately from cross-system compatibility.

\subsection{Rank-changing defect states}
A defect can introduce or remove a localized state from the selected window. Three
scientifically distinct responses are permitted:
\begin{enumerate}
 \item revise both retained spaces under a reviewed physical rule and repeat every
 projection test;
 \item choose rank-$q$ projectors $Q_b\le P_b$ and $Q_d\le P_d$ and construct only the
 explicitly named common-core perturbation $\Delta H_d^{(q)}$, retaining omitted-space
 weights and couplings as error channels; or
 \item block subtraction when neither construction is physically justified.
\end{enumerate}
Zero padding, duplicate orbitals, or deletion chosen after inspecting the desired answer
do not create a common space. A separate defect-only sector may be retained in a block
operator, but host-to-defect subtraction is then defined only on the admitted common
block and its couplings to the extra sector remain visible.

\subsection{Controlled construction with guaranteed properties}
A bounded synthetic benchmark can construct every premise of the common-space operation
before asking numerical code to recover it. Let the finite ambient dimension satisfy
$N\geq2M$, and choose periodic $N\times M$ frames $F_b$ and $E_b$ such that
$[F_b\ E_b]$ has orthonormal columns. Choose a periodic $G(\mathbf k)\in U(M)$ and
principal-angle fields
\begin{equation}
 0\leq\theta_j(\mathbf k)\leq\theta_{\max}<\frac{\pi}{2},
 \qquad
 C=\operatorname{diag}(\cos\theta_j),
 \qquad
 S=\operatorname{diag}(\sin\theta_j),
\end{equation}
and define
\begin{equation}
 F_d=F_bG^\dagger C+E_bS,
 \qquad
 P_s=F_sF_s^\dagger.
 \label{eq:constructed-defect-frame}
\end{equation}
These definitions give exact identities
\begin{equation}
 F_d^\dagger F_d=I_M,
 \qquad
 F_d^\dagger F_b=CG,
 \qquad
 \sigma_j(F_d^\dagger F_b)=\cos\theta_j,
 \label{eq:constructed-overlap}
\end{equation}
so the cross-overlap is full rank and its unitary polar factor is exactly $G$. Thus
\begin{equation}
 U=F_dGF_b^\dagger
 \label{eq:constructed-alignment}
\end{equation}
is a supplied periodic partial unitary from $\operatorname{Range}P_b$ onto
$\operatorname{Range}P_d$. Rank, periodic sewing, principal angles, conditioning, and
the intended alignment are construction guarantees rather than conclusions inferred
from a sampled optimizer.

To construct parent operators as well, choose periodic Hermitian retained fields $K_s$,
scalar references $e_s$, and retained-relative exterior levels $\Lambda_s$ satisfying
\begin{equation}
 \Lambda_s-\lambda_{\max}(K_s(\mathbf k))\geq\Delta_s>0,
\end{equation}
and set
\begin{equation}
 H_s=F_sK_sF_s^\dagger+\Lambda_s(I-P_s)+e_sI_N.
 \label{eq:constructed-parent-pair}
\end{equation}
After subtracting the supplied scalar references, compression, transport, and subtraction
then obey
\begin{equation}
 \overline H_s^{(P)}=F_sK_sF_s^\dagger,
\end{equation}
\begin{equation}
 \Delta H_d^{(P)}
 =F_d\left(K_d-GK_bG^\dagger\right)F_d^\dagger.
 \label{eq:constructed-defect-difference}
\end{equation}

The grammar exposes compatible controls separately: $M$ fixes retained rank;
$\theta_j$ fixes overlap conditioning; $G$ fixes the admitted coordinate map; $K_b$ and
$K_d$ fix the aligned retained operators; $e_s$ tests scalar-reference recovery;
$\Lambda_s$ and $\Delta_s$ fix exterior isolation; and the harmonic content of
$F_b,E_b,G,$ and $K_s$ influences real-space support. Exact construction identities can
be checked symbolically or by an independent formula route. Numerical work measures
what the grammar does not fix, including recovery error under finite precision, response
to controlled perturbations, real-space hopping tails, truncation error, and gate
sensitivity.

This is not a construction of an arbitrary physical defect. The global periodic frames
make both retained bundles complex-trivial; semantic site and orbital correspondence is
stipulated by the benchmark labels rather than discovered; and not every desired
locality, symmetry, spectrum, and principal-angle profile can be imposed independently.
Rank-changing, topologically obstructed, or materially derived cases require the separate
routes already stated. Paper 3~\cite{Ragasa2026c} supplies the higher-rank gauge-circuit
grammar from which $F_b$, $E_b$, and $G$ may be drawn. Paper 5 adds two-system
principal-angle, energy-zero, and subtraction controls. Paper 6 replaces global frames
by patchwise data when topology
or a required symmetry category is itself under test. This relationship transfers no
mesh, threshold, numerical result, or evidence status between papers.

\section{Alignment map}

\subsection{Cross-overlap and principal angles}
Let $F_s(\mathbf k):\mathbb C^M\rightarrow\mathcal H_{\mathrm{amb},\mathbf k}$ be
orthonormal frame matrices. The cross-overlap and its singular-value decomposition are
\begin{equation}
 S_d(\mathbf k)=F_d(\mathbf k)^\dagger F_b(\mathbf k)
 =L_d(\mathbf k)\Sigma_d(\mathbf k)R_d(\mathbf k)^\dagger.
\end{equation}
The singular values obey $\sigma_j=\cos\theta_j$ and define principal angles between the
retained subspaces. Full rank is required for the polar candidate
\begin{equation}
 G_d^{\mathrm{pol}}(\mathbf k)=L_d(\mathbf k)R_d(\mathbf k)^\dagger,
\end{equation}
which maps pristine frame coordinates to defect frame coordinates. The associated fiber
map is
\begin{equation}
 U_d(\mathbf k)=F_d(\mathbf k)G_d(\mathbf k)F_b(\mathbf k)^\dagger
 :\mathcal E_{b,\mathbf k}\rightarrow\mathcal E_{d,\mathbf k}.
\end{equation}
A small $\sigma_{\min}$ indicates an ill-conditioned physical identification even when
the polar factor is exactly unitary. Thresholds for rank, conditioning, principal
angles, and mesh refinement must be frozen before execution.

Overlap maximization is meaningful only in a common ambient representation. It is not a
physical labeling rule by itself. Degenerate singular values can leave equally optimal
polar correspondences that must be resolved, or retained as alignment uncertainty, using
independent semantics.

\subsection{Global periodicity and topology screen}
A pointwise map is admitted globally only when $U_d(\mathbf k)$ is periodic, sufficiently
regular for the intended representation, and compatible with reciprocal-face sewing.
Equal rank does not guarantee this. Paper 5 applies only a minimal screen:
\begin{itemize}
 \item verify constant rank and projector continuity;
 \item compute declared seam and loop residuals on refined meshes;
 \item attempt an explicit periodic alignment and test its closure; and
 \item escalate suspected complex, time-reversal, crystalline, or real-bundle
 obstruction to Paper 6 rather than assigning a topology label here.
\end{itemize}
A nonidentity Wilson loop is not by itself a failure. Holonomy can be nontrivial in a
bundle that still admits a global frame. A small finite-grid seam residual is not a
continuum triviality certificate.

\subsection{Orbital semantics and constrained correspondence}
In a localized representation, the common coordinate label is
\begin{equation}
 a=(\mathbf R,\tau,\mu,\eta),
\end{equation}
where $\mathbf R$ is a comparison-lattice translation, $\tau$ a site, $\mu$ an orbital
semantic, and $\eta$ an internal degree of freedom. A candidate correspondence must
satisfy frozen constraints on translations, centers, sites, orbital character, internal
coordinate order, interpolation continuity when used, and representations of symmetries
that survive in both systems.

The defect generally lowers symmetry. Only the common surviving subgroup may constrain
the map. If $D_s(g,\mathbf k)$ represents a required unitary operation $g$, an equivariant
map obeys
\begin{equation}
 D_d(g,\mathbf k)U_d(\mathbf k)
 =U_d(g\mathbf k)D_b(g,\mathbf k).
\end{equation}
Antiunitary operations require their explicit conjugation rule rather than this unitary
shorthand.

Assignment may be formulated as constrained matching, but an optimizer score is not a
scientific certificate. The admissible set, tie handling, assignment margins, and
sensitivity to windows and localization starts are retained. If physically different
pairings remain indistinguishable, the alignment is ambiguous.

\subsection{Admission diagnostics}
At minimum, a future protocol freezes norms and tolerances for
\begin{equation}
 \varepsilon_{\mathrm{unit}}
 =\max_{\mathbf k}\lVert G_d^\dagger G_d-I\rVert,
 \qquad
 \sigma_{\min}=\min_{\mathbf k,j}\sigma_j(\mathbf k),
\end{equation}
\begin{equation}
 \varepsilon_{\mathrm{sew}}
 =\max_f\lVert U_d(\mathbf k_f^+)
 -T_d(f)U_d(\mathbf k_f^-)T_b(f)^\dagger\rVert,
\end{equation}
and applicable symmetry, center, character, assignment-margin, and mesh-refinement
residuals. $T_s(f)$ is the declared reciprocal-face sewing map. These channels are not
combined into one score. Failure of any mandatory channel blocks the comparison.

\section{Energy alignment and impurity construction}

\subsection{Scalar energy reference}
Independent Kohn--Sham operators are defined only up to an additive scalar for this
comparison. Let
\begin{equation}
 \overline H_s^{(P)}(\mathbf k)
 =H_s^{(P)}(\mathbf k)-E_{\mathrm{ref},s}I.
\end{equation}
The same physical estimator must be used for both branches: for example, an approved
bulk-like electrostatic region, compatible deep state, or another reproducible
reference. The estimator, region, finite-supercell drift, parameter sensitivity, and
uncertainty are retained. A band edge may be used only if the same state identity and
finite-size convention are maintained.

Energy alignment removes one scalar identity component. It does not correct an
incompatible projector, gauge, lattice, orbital assignment, spinor order, or local
electrostatic response.

\subsection{Transport and subtraction}
The pristine operator transported to the defect fiber is
\begin{equation}
 H_{b\rightarrow d}^{(P)}(\mathbf k)
 =U_d(\mathbf k)\overline H_b^{(P)}(\mathbf k)U_d(\mathbf k)^\dagger.
\end{equation}
The aligned defect-induced operator is
\begin{equation}
 \boxed{\Delta H_d^{(P)}(\mathbf k)
 =\overline H_d^{(P)}(\mathbf k)-H_{b\rightarrow d}^{(P)}(\mathbf k)}.
\end{equation}
In aligned Wannier coordinates,
\begin{equation}
 \Delta H_{W,d}(\mathbf R)
 =\frac{1}{|\mathbb T_c^D|}\int_{\mathbb T_c^D}
 e^{-i\mathbf k\cdot\mathbf R}\Delta H_d^{(P)}(\mathbf k)\,d\mathbf k,
\end{equation}
with a declared discrete-transform convention in computation. Hermiticity is checked as
\begin{equation}
 \Delta H_{W,d}(-\mathbf R)=\Delta H_{W,d}(\mathbf R)^\dagger.
\end{equation}

Direct matrix subtraction is permitted only after the coordinates have been aligned so
that the matrix of $U_d$ is the identity. Equal shape alone is insufficient.

\subsection{Gauge covariance}
Under independent retained-space frame changes $V_b$ and $V_d$,
\begin{equation}
 H_s\mapsto V_s^\dagger H_sV_s,
 \qquad G_d\mapsto V_d^\dagger G_dV_b,
\end{equation}
and therefore
\begin{equation}
 \Delta H_d\mapsto V_d^\dagger\Delta H_dV_d.
\end{equation}
The spectrum, singular values, traces, and unitarily invariant norms are unchanged.
Individual onsite and hopping entries, orbital decompositions, and apparent spatial
support generally are not. They are reportable only in the frozen aligned localized
gauge.

A controlled covariance attack recomputes the full pipeline after independent admissible
frame rotations and checks
\begin{equation}
 \varepsilon_{\mathrm{cov}}
 =\lVert\Delta H_d'-V_d^\dagger\Delta H_dV_d\rVert.
\end{equation}
Rotating only the final matrix tests algebra, not the alignment implementation.

\subsection{Independent reconstruction}
A second implementation must reconstruct the transported pristine operator and defect
difference without importing the primary alignment routine. It verifies spectra,
Hermiticity, trace identities, selected real-space blocks, covariance, and reciprocal/
real-space reconstruction. Agreement establishes implementation consistency under the
frozen contract, not physical validation of the projector or parent model.

\section{Validation, failure semantics, and specialization}

\subsection{Separated error ledger}
The protocol retains parent-model; basis, cutoff, and Brillouin-mesh; finite-supercell;
folding; projection and disentanglement; localization and optimizer-basin; common-space
alignment; scalar-reference; Fourier-transform and aliasing; and later impurity-reduction
errors as distinct channels. They are not quadratically combined without a justified
statistical model. A relative error with zero reference norm is replaced by its declared
absolute form.

\subsection{Fail-closed dispositions}
\begin{center}\small
\begin{tabularx}{\textwidth}{>{\bfseries\raggedright\arraybackslash}p{0.32\textwidth}Y}
\toprule
Disposition & Meaning \\
\midrule
\nolinkurl{COMMON_SPACE_ADMITTED} & All mandatory geometry, rank, global, semantic, energy, and verification gates pass. \\
\nolinkurl{COMMON_CORE_ONLY} & A reviewed lower-rank common block is admitted; omitted states and couplings remain explicit. \\
\nolinkurl{RANK_INCOMPATIBLE} & Retained content differs and no approved common core or enlargement exists. \\
\nolinkurl{ALIGNMENT_ILL_CONDITIONED} & Cross-overlap is singular or below its conditioning gate. \\
\nolinkurl{GLOBAL_COMPATIBILITY_UNRESOLVED} & Periodic or topological compatibility is not certified. \\
\nolinkurl{SEMANTIC_ASSIGNMENT_AMBIGUOUS} & More than one physically distinct correspondence survives. \\
\nolinkurl{ENERGY_ALIGNMENT_UNRESOLVED} & No stable common scalar reference passes its gate. \\
\nolinkurl{NUMERICALLY_UNRESOLVED} & Refinement or independent reconstruction fails without diagnosing structural obstruction. \\
\bottomrule
\end{tabularx}
\end{center}
A blocked disposition does not create a zero impurity operator. Optional human choices,
such as revising the projector family, remain separate from technical review and require
a new versioned contract.

\subsection{Silicon specialization contracts}
A later P:Si or B:Si specialization must freeze separate pristine/defect pairs,
pseudopotentials, charge and magnetic states, structural rules, projectors, spin
conventions, supercells, alignment regions, thresholds, and finite-size sequences. The
two dopants do not inherit one another's orbital assignment or tolerances.

Paper 4's scalar-relativistic non-SOC ten-spatial-orbital model is a predecessor, not the
direct common-space parent for an SOC defect calculation. Paper 7 therefore requires
matched pristine-Si/SOC references. Its proposed SOC P:Si branch is also blocked by the
current physical specification, which selects a scalar-relativistic non-SOC primary
P:Si branch; only a reviewed amendment or successor can change that selection.
\end{prospective}

\section{Current evidence boundary}
No maintained experiment under \texttt{experiments/} implements this common-space defect
construction. No matched pristine/defect calculation, ambient cross-overlap, admitted
alignment map, energy-alignment estimate, common-core construction, extracted impurity
operator, independent verifier, or immutable Paper 5 evidence package exists.

Synthetic gauge and locality calculations associated with Papers 1--3 can motivate
software diagnostics but do not instantiate a pristine/defect comparison. Paper 4 has no
production bulk-silicon parent, and Paper 7 has no production dopant branch. Historical
files under \texttt{calculations/} do not transfer their identities, tolerances, or
evidence status into this paper.

Accordingly, this manuscript supports definitions, decision gates, and a prospective
record schema only. It does not establish that P:Si, B:Si, or another material pair has
compatible retained spaces. It reports no impurity matrix element, range, spectrum,
binding energy, or continuum parameter.

\section{Discussion}
The construction separates three questions: do the calculations use compatible periodic
geometry; do their retained projector families admit the required global identification;
and does a particular alignment respect the scientific semantics assigned to its
coordinates? A positive answer to one does not imply the others.

The polar factor is useful precisely because its role is limited. It gives the nearest
unitary to a full-rank cross-overlap, but it cannot decide whether two orbitals represent
the same site or physical channel. Topology is likewise relative to the required global
structure. Two individually nontrivial bundles can still be relatively isomorphic, while
a symmetry-compatible identification can be obstructed even when an unconstrained
complex unitary exists. Paper 6 owns those distinctions.

A common-core construction can preserve useful information in a rank-changing problem,
but it changes the scientific object. Its result is a projected common-core perturbation,
not the full defect operator. Discarded weights and couplings are part of the result.

Gauge covariance is not gauge independence of entries. Local impurity language such as
``onsite shift'' or ``modified hopping'' becomes meaningful only after a localized and
semantically aligned gauge has been declared. Operator invariants can test the
construction; they cannot make an arbitrary decomposition unique.

\section{Conclusions}
A periodic defect operator is obtained by transport and subtraction on a certified
common retained space. Required evidence includes matched geometry, compatible
projectors, a regular periodic alignment, physical orbital correspondence, stable scalar
energy alignment, covariance, and independent reconstruction. Rank change, global
obstruction, semantic ambiguity, and numerical nonconvergence receive distinct
fail-closed dispositions.

The framework is ready to guide reviewed experiment design, but no immutable Paper 5
protocol or execution exists. Before material work, the project must select a bounded
synthetic common-space benchmark, schemas and verifier identities, exact tolerances, and
then separate P:Si and B:Si specializations consistent with the physical specification.

\appendix
\section{Pre-execution commitment ledger}
\begin{center}\small
\begin{tabularx}{\textwidth}{p{0.22\textwidth}YY}
\toprule
\textbf{Item} & \textbf{Required commitment} & \textbf{Current state} \\
\midrule
Comparison pair & Exact pristine/defect identities and physical branches & Unselected \\
Geometry & Common lattice, folding route, site and translation maps & Method only \\
Parent numerics & Code, pseudopotentials, cutoffs, meshes, tolerances & Unselected \\
Retained spaces & Ranks, windows, projectors, defect-state policy & Unselected \\
Global structure & Sewing conventions and minimal topology screen & Method only \\
Correspondence & Center, site, orbital, internal, symmetry constraints & Method only \\
Alignment gates & Norms, singular-value, seam, semantic, refinement thresholds & Unselected \\
Energy reference & Estimator, region, uncertainty, stability rule & Unselected \\
Verification & Independent route, attacks, reconstruction tolerances & Unselected \\
Evidence & Immutable version, schemas, package and verifier identities & Absent \\
Resources & Executables, inputs, scale, outputs, runtime and memory & Unreported \\
\bottomrule
\end{tabularx}
\end{center}

\section{Required result record}
A future retained Paper 5 result must contain:
\begin{enumerate}
 \item immutable parent, code, environment, and pseudopotential identities;
 \item pristine and defect geometry plus folding and translation maps;
 \item projector ranks, windows, spectral gaps, and defect-state dispositions;
 \item all cross-overlap singular values and principal-angle summaries;
 \item periodic-sewing, topology-screen, and surviving-symmetry residuals;
 \item orbital assignments, constraints, margins, ties, and localization sensitivity;
 \item the alignment map and its unitarity, conditioning, and covariance diagnostics;
 \item energy-reference estimates, regions, sensitivities, and uncertainties;
 \item complete reciprocal- and real-space impurity operators with Hermiticity checks;
 \item independent reconstruction and controlled frame-attack results;
 \item separated numerical, finite-size, projection, localization, alignment, and transform errors;
 \item every admitted, failed, blocked, and unresolved disposition; and
 \item manifests, hashes, commands, platform, schema, run, verifier, and evidence-package identities.
\end{enumerate}

\bibliographystyle{unsrt}
\bibliography{references,../../unpublished-papers}

\end{document}
